\documentclass[11pt]{article}
\usepackage{amsmath,amssymb,amsthm}
\usepackage[margin=1.1in]{geometry}
\usepackage{booktabs}
\usepackage{hyperref}

\newtheorem{proposition}{Proposition}
\newtheorem{remark}{Remark}

\title{Byte-Pair Encoding:\\
Designing the Alphabet, and Paying for It Honestly}
\author{Yuval Lavie}
\date{\today}

\begin{document}
\maketitle

\begin{abstract}
Every model on the ladder read text one character at a time --- a
choice made by default, never examined. This note examines it. We
implement BPE by hand (greedy pair-merging over a pre-tokenized
corpus), verify it is a lossless bijection, watch it rediscover
English units by pure counting, and then rerun rung 4's GPT unchanged
on the new 515-token alphabet. The comparison is made in the only
currency that survives a change of tokenizer --- NLL \emph{per
character} --- and the result is a genuine experiment: the BPE model
overfits on schedule (the token corpus is $2.3\times$ smaller),
early stopping enters the toolkit, and the character alphabet wins at
this scale. \texttt{bpe\_theoretical.py},
\texttt{bpe\_practical\_pytorch.py}.
\end{abstract}

\section{The algorithm, and what it learns}

A tokenizer is a bijection between text and sequences over a
vocabulary we design. BPE builds the vocabulary greedily from data:
start from the characters; repeatedly merge the most frequent adjacent
pair into a new token. Each merge is the single vocabulary addition
that maximizes the immediate drop in corpus length --- greedy
compression. Following GPT-2, merges act only inside \emph{chunks} (a
word with its preceding whitespace), which keeps tokens interpretable
and the training loop cheap (unique chunks weighted by count, not the
raw stream).

Run on Tiny Shakespeare's training split ($450$ merges): the first
merges are \texttt{'\ t'}, \texttt{'he'}, \texttt{'ou'},
\texttt{'\ the'}, \texttt{'\textbackslash n\textbackslash n'} (the
blank line between speeches!); by the end the vocabulary holds
\texttt{'\ should'}, \texttt{'\ father'}, \texttt{'ICHARD'} (from the
dramatis personae). English structure, discovered by counting ---
rung 2's embedding lesson, one level down the stack.

Verified: encoding round-trips the full $1{,}115{,}394$-character
corpus \emph{exactly} (a tokenizer bug is silent data corruption; the
first implementation failed this assert by dropping trailing
whitespace, and the fix is documented in the code); compression is
monotone in vocabulary size, reaching $2.33$ chars/token at $V = 515$.

\section{The measuring-stick trap}

Perplexity per \emph{token} is meaningless across tokenizers: a model
over bigger tokens takes fewer, harder guesses. The exchange rate is
the compression ratio, so all comparisons are made per character:
\begin{equation}
  \boxed{\;\mathrm{NLL/char}
  \;=\; \frac{\sum_i \mathrm{NLL}(\mathrm{token}_i)}
             {\sum_i |\mathrm{token}_i|}\;}
  \label{eq:perchar}
\end{equation}
(total validation NLL divided by the characters the evaluated tokens
cover). Hygiene: the merges are trained on the \emph{training split
only} --- a tokenizer fitted on validation text would smuggle val
statistics into the alphabet.

\section{The experiment: same GPT, new alphabet}

Rung 4's architecture, unchanged, with $V = 515$: a 128-token window
now spans $\approx 298$ characters ($2.3\times$ the context), and
\texttt{'\ the'} is one prediction instead of four. The opposing
force: the training corpus shrinks from $10^6$ char-positions to
$431{,}399$ token-positions.

The opposing force wins the late game. Rung 4's 6000-step budget is
$\sim$113 epochs over the token corpus; validation NLL/char bottoms
at step $1250$ and then climbs a quarter nat while train loss keeps
falling --- overfitting, arriving $2.3\times$ sooner exactly as the
data shrinkage predicts. \textbf{Early stopping} (keep the best
validation checkpoint, not the last) enters the toolkit here, at the
first place its absence measurably costs:
\begin{center}
\begin{tabular}{lrr}
\toprule
model & val NLL/char & bits/char\\
\midrule
BPE-GPT, last step (overfit) & $1.7904$ & $2.583$\\
BPE-GPT, early-stopped (step 1250) & $1.5165$ & $2.188$\\
char-GPT (rung 4) & $\mathbf{1.4724}$ & $\mathbf{2.124}$\\
\bottomrule
\end{tabular}
\end{center}

\begin{remark}[The verdict, and its scope]
At this scale \emph{characters win}: when data --- not context length
--- is the bottleneck, a bigger alphabet just makes the dataset
smaller. BPE's dominance in real LLMs lives at the opposite corner:
corpora of $10^{12}$ tokens where data is plentiful, context is the
scarce resource, and spelling is a waste of depth. The point of this
note is not the winner but the method: \eqref{eq:perchar} turned
``which tokenizer is better?'' into an answerable question.
\end{remark}

\begin{remark}[Tokenizer lock-in]
The BPE model's checkpoint is only meaningful together with its merge
list: the tokenizer is part of the model. This constraint returns in
the finetuning note, where a new corpus must be expressed in the old
model's alphabet.
\end{remark}

\end{document}
